\chapter{Exponentiële functie en logaritme}\label{ch:g12:exp}

De \omterm{thm:g12:exp:existence}{exponentiële functie} is de unieke \omterm{def:g11:func:function}{functie} die gelijk is aan haar eigen
\omterm{def:g12:deriv:derivative}{afgeleide} en in $0$ de waarde $1$ aanneemt. Ze zet sommen om in producten;
haar inverse, de \omterm{def:g12:exp:ln}{natuurlijke logaritme}, zet producten om in sommen. Samen
beschrijven ze elk verschijnsel waarvan de veranderingssnelheid evenredig is
met de omvang: radioactief verval, bevolkingsgroei, samengestelde intrest.

\section{De exponentiële functie}

\begin{theorem}[Bestaan en uniciteit]\label{thm:g12:exp:existence}
Er bestaat precies één differentieerbare \omterm{def:g11:func:function}{functie} $f \colon \R \to \R$ waarvoor
\[
f' = f \qquad\text{en}\qquad f(0) = 1 .
\]
Ze heet de \emph{exponentiële functie}\index{exponentiële functie} en wordt
$\exp$ geschreven, of $x \mapsto \eu^x$.
\end{theorem}

\begin{proof}[Bewijs van de uniciteit]
Zo'n \omterm{def:g11:func:function}{functie} verdwijnt om te beginnen nooit. Neem immers
$g(x) = f(x)f(-x)$; dan is
\[
g'(x) = f'(x)f(-x) - f(x)f'(-x) = f(x)f(-x) - f(x)f(-x) = 0,
\]
dus is $g$ constant en gelijk aan $g(0) = 1$: voor elke $x$ is
$f(x)f(-x) = 1$, en in het bijzonder $f(x) \neq 0$.

Zij nu $f_1, f_2$ twee oplossingen en $h = \dfrac{f_1}{f_2}$ (geoorloofd,
want $f_2$ verdwijnt nooit). Dan is
\[
h' = \frac{f_1' f_2 - f_1 f_2'}{f_2^2} = \frac{f_1 f_2 - f_1 f_2}{f_2^2} = 0,
\]
zodat $h$ constant is en gelijk aan $h(0) = 1$, dus $f_1 = f_2$.

Het \emph{bestaan} wordt op dit niveau aangenomen (het volgt uit de \omterm{def:g12:seq:sequence}{rij} van
\cref{exo:g12:seq:10}, of uit de inverse van de \omterm{def:g12:exp:ln}{logaritme} die in
\cref{ch:g12:integ} met integratie opgebouwd wordt).
\end{proof}

\begin{proposition}[Functionaalvergelijking]\label{prop:g12:exp:funceq}
Voor alle $x, y \in \R$ en $n \in \Z$ geldt
\[
\eu^{x+y} = \eu^x\,\eu^y, \qquad
\eu^{-x} = \frac{1}{\eu^x}, \qquad
\eu^{x-y} = \frac{\eu^x}{\eu^y}, \qquad
\bigl(\eu^{x}\bigr)^n = \eu^{nx}.
\]
Bovendien is $\eu^x > 0$ voor alle $x$.
\end{proposition}

\begin{proof}
Kies $y$ vast en bekijk $\varphi(x) = \dfrac{\exp(x+y)}{\exp(x)\exp(y)}$. De
teller en de noemer zijn, als \omterm{def:g11:func:function}{functies} van $x$, allebei oplossingen van
$f' = f$ op de constante $\exp(y)$ na; $\varphi$ rechtstreeks afleiden
(quotiëntregel) geeft $\varphi' = 0$, dus $\varphi \equiv \varphi(0) = 1$,
wat de eerste identiteit bewijst. Met $y = -x$ krijg je de tweede (met de
waarde $\eu^0 = 1$), en de derde volgt daaruit. De vierde is voor $n \geq 0$
een inductie op de eerste, en voor $n < 0$ een gevolg van de tweede.

Positiviteit: $\eu^x = \left(\eu^{x/2}\right)^2 \geq 0$ en $\eu^x \neq 0$
(aangetoond in \cref{thm:g12:exp:existence}), dus $\eu^x > 0$.
\end{proof}

\begin{proposition}[Verloop en limieten]\label{prop:g12:exp:variations}
De \omterm{thm:g12:exp:existence}{exponentiële functie} is strikt \omterm{def:g12:seq:monotonic}{stijgend} en \omterm{def:g12:deriv:convex}{convex}, en
\[
\lim_{x\to-\infty} \eu^x = 0, \qquad
\lim_{x\to+\infty} \eu^x = +\infty .
\]
\end{proposition}

\begin{proof}
$(\exp)' = \exp > 0$ geeft de strikte stijging; $(\exp)'' = \exp > 0$ geeft
de \omterm{def:g12:deriv:convex}{convexiteit}. Door de \omterm{def:g12:deriv:convex}{convexiteit} is $\eu^x \geq 1 + x$ (\omterm{def:g12:deriv:tangent}{raaklijn} in $0$,
zie \cref{ex:g12:deriv:ineq}), dus $\eu^x \to +\infty$ voor $x \to +\infty$
door \omterm{def:g10:algebra:equation}{vergelijking}. Daarna is $\eu^{x} = 1/\eu^{-x} \to 0$ voor
$x \to -\infty$.
\end{proof}

\begin{omfigure}
\begin{tikzpicture}
\begin{axis}[omaxis,
  width=8.8cm, height=6cm,
  xmin=-3.3, xmax=2.4, ymin=-2.4, ymax=7.8,
  xtick={-2,-1,1,2}, ytick={1}]
  \addplot[omProp, thick, domain=-3.2:2.3] {1 + x};
  \addplot[omDef, thick, domain=-3.2:2.0] {exp(x)};
  \fill (axis cs:0,1) circle (1.5pt);
  \node[omDef, font=\small, anchor=east] at (axis cs:1.35,5.2) {$\eu^x$};
  \node[omProp, font=\small, anchor=north west] at (axis cs:1.05,1.6) {$y = 1+x$};
\end{axis}
\end{tikzpicture}

{\small De \omterm{thm:g12:exp:existence}{exponentiële functie}: strikt \omterm{def:g12:seq:monotonic}{stijgend}, \omterm{def:g12:deriv:convex}{convex}, met
$\lim_{-\infty} \eu^x = 0$. Ze ligt boven haar \omterm{def:g12:deriv:tangent}{raaklijn} in $0$:
$\eu^x \geq 1 + x$.}
\end{omfigure}

\begin{theorem}[Vergelijking van groeiorden]\label{thm:g12:exp:growth}
\index{groeiorden}
Voor elk \omterm{def:g10:numbers:sets}{geheel getal} $n \geq 1$ geldt
\[
\lim_{x\to+\infty} \frac{\eu^x}{x^n} = +\infty,
\qquad
\lim_{x\to-\infty} x^n \eu^x = 0 .
\]
In woorden: de \omterm{thm:g12:exp:existence}{exponentiële functie} verslaat elke macht van $x$.
\end{theorem}

\begin{proof}
Voor $n = 1$: pas $\eu^t \geq 1 + t \geq t$ toe met $t = x/2$:
\[
\frac{\eu^x}{x} = \frac{\left(\eu^{x/2}\right)^2}{x}
\geq \frac{(x/2)^2}{x} = \frac{x}{4} \xrightarrow[x\to+\infty]{} +\infty .
\]
Voor algemene $n$ schrijf je
\[
\frac{\eu^x}{x^n}
= \frac{1}{n^n}\left(\frac{\eu^{x/n}}{x/n}\right)^{n}
\xrightarrow[x\to+\infty]{} +\infty
\]
met het geval $n=1$ en de samenstelling. De limiet in $-\infty$ volgt met de
substitutie $x \mapsto -x$:
$\abs{x^n \eu^x} = \frac{\abs{x}^n}{\eu^{\abs x}} \to 0$.
\end{proof}

\begin{omfigure}
\begin{tikzpicture}
\begin{axis}[omaxis,
  width=8.8cm, height=6cm,
  xmin=-0.4, xmax=5.6, ymin=-12, ymax=155,
  xtick={1,2,3,4,5}, ytick={50,100,150}]
  \addplot[omThm, thick, domain=0:5.4] {x^2};
  \addplot[omProp, thick, domain=0:5.3] {x^3};
  \addplot[omDef, thick, domain=-0.4:5.0] {exp(x)};
  \fill (axis cs:4.536,93.4) circle (1.5pt);
  \draw[black, dashed, thin] (axis cs:4.536,0) -- (axis cs:4.536,93.4);
\end{axis}
\end{tikzpicture}

{\small De \omterm{thm:g12:exp:existence}{exponentiële functie} (blauw) verslaat uiteindelijk elke macht: ze
steekt $x^2$ (rood) al vroeg voorbij en $x^3$ (oranje) rond
$x \approx 4.54$.}
\end{omfigure}

\section{De natuurlijke logaritme}

\begin{definition}[Natuurlijke logaritme]\label{def:g12:exp:ln}
De \omterm{thm:g12:exp:existence}{exponentiële functie} is \omterm{def:g12:limcont:continuity}{continu} en strikt \omterm{def:g12:seq:monotonic}{stijgend} van $\R$ op
$\intoo{0}{+\infty}$; volgens de bijectiestelling
(\cref{thm:g12:limcont:bijection}) heeft de \omterm{def:g10:algebra:equation}{vergelijking} $\eu^x = y$ voor elke
$y > 0$ precies één oplossing. Die oplossing is de \emph{natuurlijke
logaritme}\index{logaritme} van $y$, geschreven $\ln y$. Dus geldt
\[
\text{voor } x \in \R,\ y > 0: \qquad y = \eu^x \iff x = \ln y .
\]
In het bijzonder is $\ln 1 = 0$, $\ln \eu = 1$, $\eu^{\ln y} = y$ en
$\ln(\eu^x) = x$.
\end{definition}

\begin{omfigure}
\begin{tikzpicture}
\begin{axis}[omaxis,
  width=8.6cm, axis equal image,
  xmin=-4.3, xmax=6.3, ymin=-4.3, ymax=6.3,
  xtick={1}, ytick={1}]
  \addplot[black, dashed, domain=-3.9:5.9] {x};
  \addplot[omDef, thick, domain=-4.1:1.79] {exp(x)};
  \addplot[omThm, thick, domain=0.0167:5.9, samples=120] {ln(x)};
  \node[omDef, font=\small, anchor=east] at (axis cs:0.95,2.6) {$\eu^x$};
  \node[omThm, font=\small, anchor=north] at (axis cs:4.3,1.25) {$\ln x$};
  \node[font=\small, anchor=west, rotate=45] at (axis cs:4.6,4.85) {$y=x$};
\end{axis}
\end{tikzpicture}

{\small De krommen van $\exp$ en $\ln$ zijn elkaars spiegelbeeld om de rechte
$y = x$: het punt $(x, \eu^x)$ spiegelt naar $(\eu^x, x)$.}
\end{omfigure}

\begin{proposition}[Algebraïsche eigenschappen]\label{prop:g12:exp:lnalg}
Voor alle $a, b > 0$ en $n \in \Z$ geldt
\[
\ln(ab) = \ln a + \ln b, \quad
\ln\frac{1}{a} = -\ln a, \quad
\ln\frac{a}{b} = \ln a - \ln b, \quad
\ln(a^n) = n \ln a, \quad
\ln\sqrt{a} = \tfrac12 \ln a .
\]
\end{proposition}

\begin{proof}
$\eu^{\ln a + \ln b} = \eu^{\ln a}\eu^{\ln b} = ab$, en van beide leden de
\omterm{def:g12:exp:ln}{logaritme} nemen geeft de eerste identiteit. De andere volgen met hetzelfde
mechanisme uit de overeenkomstige identiteiten van
\cref{prop:g12:exp:funceq}.
\end{proof}

\begin{proposition}[Analytische eigenschappen]\label{prop:g12:exp:lnanalytic}
De \omterm{def:g11:func:function}{functie} $\ln$ is \omterm{def:g12:deriv:derivative}{differentieerbaar} op $\intoo{0}{+\infty}$ met
\[
(\ln x)' = \frac{1}{x},
\]
strikt \omterm{def:g12:seq:monotonic}{stijgend}, \omterm{def:g12:deriv:convex}{concaaf}, en $\lim\limits_{x\to0^+} \ln x = -\infty$,
$\lim\limits_{x\to+\infty} \ln x = +\infty$. Bovendien geldt voor elk \omterm{def:g10:numbers:sets}{geheel
getal} $n \geq 1$:
\[
\lim_{x\to+\infty} \frac{\ln x}{x^{1/n}} = 0, \qquad
\lim_{x\to0^+} x \ln x = 0 .
\]
\end{proposition}

\begin{proof}
De differentieerbaarheid van de inverse \omterm{def:g11:func:function}{functie} wordt op dit niveau
aangenomen; ga je ervan uit, leid dan de identiteit $\eu^{\ln x} = x$ af met
de kettingregel: $(\ln x)'\,\eu^{\ln x} = 1$, dus $(\ln x)' = \frac{1}{x} > 0$,
waaruit de strikte stijging, en $(\ln x)'' = -\frac1{x^2} < 0$, waaruit de
concaviteit. De limieten in $0^+$ en $+\infty$ spiegelen die van $\exp$ door
de bijectie. Voor de vergelijking van \omterm{thm:g12:exp:growth}{groeiorden} substitueer je
$x = \eu^{t}$: $\frac{\ln x}{x} = \frac{t}{\eu^t} \to 0$ voor
$t \to +\infty$ volgens \cref{thm:g12:exp:growth}, en net zo voor de andere
limieten met $x = \eu^{-t}$, wat $x\ln x = -t\eu^{-t} \to 0$ geeft.
\end{proof}

\begin{method}[Vergelijkingen met $\exp$ en $\ln$ oplossen]\label{met:g12:exp:equations}
Beide \omterm{def:g11:func:function}{functies} zijn strikt \omterm{def:g12:seq:monotonic}{stijgend}, zodat je ze op beide leden van een
\omterm{def:g10:algebra:equation}{vergelijking} of ongelijkheid mag toepassen (of eraf halen) zonder de richting
te veranderen:
\[
\eu^{u} = \eu^{v} \iff u = v, \qquad
\eu^{u} \leq \eu^{v} \iff u \leq v,
\]
en net zo voor $\ln$ op positieve grootheden. \emph{Ga altijd eerst de
\omterm{def:g11:func:function}{domeinen} na} ($\ln$ vraagt positieve argumenten). Voor \omterm{def:g10:algebra:equation}{vergelijkingen} als
$a^x = b$ met $a>0$ en $a \neq 1$ herschrijf je $a^x = \eu^{x\ln a}$ en los je
$x = \frac{\ln b}{\ln a}$ op.
\end{method}

\begin{example}\label{ex:g12:exp:halflife}
Een radioactief monster vervalt volgens $N(t) = N_0\,\eu^{-\lambda t}$. Zijn
\emph{halveringstijd} $T$ voldoet aan $N(T) = N_0/2$, dus
$\eu^{-\lambda T} = \frac12$, waaruit $T = \frac{\ln 2}{\lambda}$: de
halveringstijd hangt niet van de beginhoeveelheid af.
\end{example}

\section{Oefeningen}

\begin{exercise}[$\star$]\label{exo:g12:exp:1}
Vereenvoudig $\dfrac{\eu^{3x}\,\eu^{-x+1}}{\eu^{x}}$ en
$\ln\!\left(\dfrac{\eu^2\sqrt{\eu}}{\eu^{-3}}\right)$.
\end{exercise}

\begin{exercise}[$\star$]\label{exo:g12:exp:2}
Los op in $\R$:
\[
\text{(a) } \eu^{2x} - 3\eu^{x} + 2 = 0;
\qquad
\text{(b) } \ln(x - 1) + \ln(x + 2) = \ln 4 .
\]
\end{exercise}

\begin{exercise}[$\star$]\label{exo:g12:exp:3}
Bereken de limieten:
\[
\lim_{x\to+\infty} \frac{\eu^x - x^2}{\eu^x + 1}, \qquad
\lim_{x\to+\infty} \frac{\ln x}{\sqrt x}, \qquad
\lim_{x\to0^+} x^2 \ln x, \qquad
\lim_{x\to+\infty} \bigl(x - \ln x\bigr).
\]
\end{exercise}

\begin{exercise}[$\star\star$]\label{exo:g12:exp:4}
Bestudeer de \omterm{def:g11:func:function}{functie} $f(x) = x\,\eu^{-x}$ op $\R$: verloop, limieten,
\omterm{def:g12:deriv:extremum}{extremum}; toon aan dat haar kromme een \omterm{def:g12:deriv:inflection}{buigpunt} heeft en geef de \omterm{def:g10:coordgeom:system}{coördinaten}
ervan.
\end{exercise}

\begin{exercise}[$\star\star$]\label{exo:g12:exp:5}
Toon aan dat voor alle $x > -1$ geldt $\ln(1 + x) \leq x$, met gelijkheid
alleen in $x = 0$. Leid af dat voor alle $n \geq 1$ geldt
\[
\left(1 + \frac{1}{n}\right)^n \leq \eu \leq \left(1 - \frac{1}{n+1}\right)^{-(n+1)} .
\]
\end{exercise}

\begin{exercise}[$\star\star$]\label{exo:g12:exp:6}
Een kapitaal $C_0$ wordt belegd tegen een jaarlijkse rentevoet van $3\%$, met
jaarlijkse kapitalisatie van de intrest.
\begin{enumerate}
  \item Druk het kapitaal $C_n$ na $n$ jaar uit.
  \item Na hoeveel jaar verdubbelt het kapitaal? Geef het exacte antwoord met
        $\ln$, en daarna een getalwaarde.
  \item Vergelijk met de \omterm{def:g10:numbers:approx}{benadering} ``$70$ gedeeld door de rentevoet in
        procent'' die bankiers gebruiken.
\end{enumerate}
\end{exercise}

\begin{exercise}[$\star\star$]\label{exo:g12:exp:7}
Los de ongelijkheid $\eu^{2x} - \eu^{x+1} > 0$ op, en daarna de ongelijkheid
$\ln(x^2 - 1) \leq \ln(x + 5)$.
\end{exercise}

\begin{exercise}[$\star\star\star$]\label{exo:g12:exp:8}
Zij $f(x) = \dfrac{\eu^x}{x}$ voor $x > 0$.
\begin{enumerate}
  \item Bestudeer het verloop van $f$ op $\intoo{0}{+\infty}$ en geef haar
        \omterm{def:g10:functions:extrema}{minimum}.
  \item Voor welke waarden van $k$ heeft de \omterm{def:g10:algebra:equation}{vergelijking} $\eu^x = kx$ in
        $\intoo{0}{+\infty}$ $0$, $1$ of $2$ oplossingen?
\end{enumerate}
\end{exercise}

\begin{exercise}[$\star\star\star$]\label{exo:g12:exp:9}
Zij voor $n \geq 1$ het getal $u_n = \left(1 + \frac1n\right)^n$.
\begin{enumerate}
  \item Toon met \cref{exo:g12:exp:5} aan dat $(u_n)$ \omterm{def:g12:seq:bounded}{naar boven begrensd} is
        door~$\eu$.
  \item Toon aan dat $\ln u_n = n \ln\left(1 + \frac1n\right) \to 1$, en leid
        af dat $u_n \to \eu$. (Tip: herken een differentiequotiënt van $\ln$
        in $1$.)
\end{enumerate}
\end{exercise}

\section{Opgave: de logaritme temt de wereld}

\begin{problem}[{Weekendopgave --- verdubbelingstijden en de regel van 72, de
schalen van aardbevingen, zuren en piano's, en de constante $\eu$ die overal
vingerafdrukken achterlaat}]
\label{pb:g12:exp:1}
Wat met een vast \emph{percentage} groeit, groeit exponentieel --- spaargeld,
bacteriën, epidemieën --- en wat te veel machten van tien overspant om te
bevatten --- energieën van aardbevingen, zuurtegraden, geluidsintensiteiten
--- wordt door een \omterm{def:g12:exp:ln}{logaritme} getemd. Deze opgave berekent
verdubbelingstijden en ontmaskert de bankiersregel van 72, leest de
logaritmische schalen van de wereld, en verzamelt de vingerafdrukken die het
getal $\eu$ op elke plaats delict achterlaat
(\cref{prop:g12:exp:funceq}, \cref{thm:g12:exp:growth},
\cref{met:g12:exp:equations}).

\medskip
\noindent\textbf{Deel I --- Vlot rekenen.}
\begin{enumerate}
  \item Los op: $\eu^{2x} = 5$; \quad $\ln(3x - 1) = 2$; \quad
        $\eu^{2x} - 3\eu^x + 2 = 0$ (een kwadratische \omterm{def:g10:algebra:equation}{vergelijking} in
        vermomming).
  \item Vereenvoudig: $\dfrac{\ln 8}{\ln 2}$; \quad
        $\ln\!\left(\eu^3 \sqrt{\eu}\right)$; \quad $\eu^{\ln 5 - \ln 2}$.
  \item Bewijs de moederongelijkheid: $\eu^x \geq 1 + x$ voor alle reële $x$
        (bestudeer $f(x) = \eu^x - 1 - x$), en leid $\ln(1 + u) \leq u$ af
        voor $u > -1$.
  \item Bereken $\lim_{x \to +\infty} x^2 \eu^{-x}$,
        $\lim_{x \to +\infty} \frac{\ln x}{x}$
        (\cref{thm:g12:exp:growth}), en
        $\lim_{x \to 0} \frac{\eu^x - 1}{x}$ (herken een \omterm{def:g12:deriv:derivative}{afgeleide}).
  \item Bestudeer $f(x) = x \ln x$ op $\intoo{0}{+\infty}$: verloop, \omterm{def:g10:functions:extrema}{minimum},
        en de limiet in $0^+$ (aangenomen: $x \ln x \to 0$). Dat kleine
        functietje meet informatie en entropie doorheen de universitaire
        volumes.
\end{enumerate}

\medskip
\noindent\textbf{Deel II --- Verdubbelingstijden en de regel van 72.}
\begin{enumerate}[resume]
  \item Spaargeld groeit met $3\,\%$ per jaar. Los $1.03^n = 2$ op: hoe lang
        duurt het voor het kapitaal verdubbelt?
  \item Bankiers schatten de verdubbelingstijd als
        $\frac{72}{\text{rentevoet in }\%}$. Toets de regel bij $3\,\%$,
        $6\,\%$ en $9\,\%$ aan de exacte $\frac{\ln 2}{\ln(1 + r)}$, en
        verklaar ze daarna: voor kleine $r$ is $\ln(1 + r) \approx r$ (de
        ongelijkheid van vraag 3 is daar de helft van het verhaal), zodat de
        exacte constante $100 \ln 2 \approx 69.3$ is --- waarom verkiezen
        bankiers $72$?
  \item Een bacterie deelt zich elke $20$ minuten: $p(t) = 2^{t/20}$ (met $t$
        in minuten). Herschrijf dat als $\eu^{\lambda t}$, en bereken daarna
        $p$ na $24$ uur. Wat bewijst het antwoord (meer dan $10^{21}$) over
        het model --- en wat houdt echte kolonies tegen?
  \item Cafeïne verlaat het lichaam met een halveringstijd van ongeveer $5$
        uur. Hoeveel blijft er na een koffie van $100$ mg om 15 uur nog over
        om 23 uur? Op welk tijdstip zakt het onder $10$ mg? (Slapeloosheid
        heeft een \omterm{def:g12:exp:ln}{logaritme}.)
  \item Eén euro tegen $100\,\%$ jaarrente: bereken het eindkapitaal bij
        jaarlijkse, maandelijkse en dagelijkse kapitalisatie, en geef het
        plafond waarvan \cref{exo:g12:exp:9} bewees dat het onbreekbaar is.
        Welke beroemde constante is de limiet van de pure hebzucht?
  \item Waarom zetten wetenschappers tijdens de beginfase van een epidemie
        $\ln(\text{aantal gevallen})$ tegen de tijd uit? Wat onthult een
        rechte lijn op die \omterm{def:g10:functions:graph}{grafiek}, en wat meet haar
        \omterm{def:g10:reffunc:affine}{richtingscoëfficiënt}?
\end{enumerate}

\medskip
\noindent\textbf{Deel III --- De logaritmische schalen van de wereld.}
(Schrijf $\log_{10} x = \frac{\ln x}{\ln 10}$.)
\begin{enumerate}[resume]
  \item Geluidsniveau in decibel: $L = 10 \log_{10}(I/I_0)$. Een gesprek meet
        $60$ dB, een rockconcert $120$ dB: met welke factor verschillen de
        geluidsintensiteiten?
  \item De magnitude van een aardbeving stijgt met $1$ wanneer de seismische
        amplitude met $10$ vermenigvuldigd wordt, en de vrijgekomen energie
        schaalt als amplitude$^{3/2}$. Vergelijk bevingen met magnitude $5$ en
        $7$: eerst de verhouding van de amplitudes, dan die van de energieën.
  \item Scheikunde: $\text{pH} = -\log_{10}[\mathrm{H^+}]$. Citroensap heeft
        pH $2$, melk pH $6.5$: wat is de verhouding van hun
        zuurconcentraties?
  \item Muziek: elk octaaf verdubbelt de \omterm{def:g10:stats:series}{frequentie}. Hoeveel octaven
        overspant een piano van haar laagste A ($27.5$ Hz) tot haar hoogste C
        ($4\,186$ Hz)? En waarom verkiezen onze zintuigen --- gehoor, zicht,
        het voelen van bevingen --- logaritmische schalen? (Eén zin.)
  \item De rekenliniaal, $350$ jaar lang de rekenmachine van de ingenieur:
        twee latjes zo geijkt dat de \emph{lengte} tot het merkteken $x$
        evenredig is met $\ln x$. Leg uit hoe het ene latje langs het andere
        schuiven getallen vermenigvuldigt, en noem de identiteit van
        \cref{prop:g12:exp:lnalg} die het werk doet.
\end{enumerate}

\medskip
\noindent\textbf{Deel IV --- De vingerafdrukken van $\eu$.}
\begin{enumerate}[resume]
  \item Uit \cref{exo:g12:exp:8}: de \omterm{def:g10:algebra:equation}{vergelijking} $\eu^x = kx$ ($x > 0$)
        heeft $0$, $1$ of $2$ oplossingen naargelang de ligging van $k$ ten
        opzichte van een drempel. Formuleer het resultaat opnieuw --- en ga
        het meetkundige feit erachter na: de rechte $y = \eu x$ is precies de
        \omterm{def:g12:deriv:tangent}{raaklijn} aan de \omterm{thm:g12:exp:existence}{exponentiële functie} door de oorsprong.
  \item De constante van het net-niet: bij $n$ loten met elk winstkans
        $\frac1n$ is $\P(\text{geen winst}) = \left(1 - \frac1n\right)^n$.
        Bereken haar limiet via $n \ln\!\left(1 - \frac1n\right)$ (gebruik de
        standaardlimiet $\frac{\ln(1 + u)}{u} \to 1$), en breng ze in
        overeenstemming met de raadselachtige $0.368$ uit
        \cref{pb:g11:binom:1}.
  \item De $37\,\%$-regel: om de beste van $n$ kandidaten te kiezen die in
        willekeurige volgorde op gesprek komen (zonder terug te kunnen),
        verwerpt de optimale strategie de eerste
        $\frac{n}{\eu} \approx 37\,\%$ en neemt ze daarna de eerste kandidaat
        die beter is dan alle vorige --- met een slaagkans van ongeveer
        $\frac1\eu$. Waar komen de twee $\frac1\eu$'s uit de vragen 18 en 19
        vandaan? Formuleer het gemeenschappelijke mechanisme (veel
        onafhankelijke gebeurtenissen met elk een kleine \omterm{def:g10:proba:distribution}{kans}), en bereken
        $\frac1\eu$ op drie decimalen.
  \item Slotstuk --- het portret van $\eu$: het plafond van de kapitalisatie
        (vraag 10); het grondtal waarvan de \omterm{def:g12:deriv:tangent}{raaklijn} in $0$ precies
        \omterm{def:g10:reffunc:affine}{richtingscoëfficiënt} $1$ heeft (vragen 3 en 4); de groei die geen
        enkele veelterm inhaalt (vraag 4); de constante van het net-niet en
        van het optimale stoppen (vragen 18--19); en haar inverse $\ln$, die
        producten in sommen omzet (vraag 16) en decaden in centimeters (deel
        III). Telkens één zin.
\end{enumerate}
\end{problem}
